\documentclass[10pt]{article}
\usepackage[margin=0.62in]{geometry}
\usepackage{booktabs}
\usepackage{graphicx}
\usepackage{microtype}
\pagestyle{plain}
\setlength{\parindent}{0pt}

\begin{document}

\begin{center}
\small
\textbf{Actual observables for the electronic-imbalance initial perturbation.}
$t_{\rm hop}=1$, $\omega_{\rm ph}=0.5$, $g/t_{\rm hop}=0.6124$
($\lambda=1.5$), $v/t_{\rm hop}=1$; wavefunction perturbation amplitude
$5\times10^{-6}$.
\end{center}

\begin{figure}[htbp]
  \centering
  \includegraphics[width=0.92\textwidth]{output/plots/paper_v_results_progression_20260804/initial_condition_raw_electronic-drive_observables.pdf}
  \caption{Actual observable trajectories for the base and perturbed initial
  states under exact cutoff-16 Hamiltonian propagation and the uncorrected
  archive moment EOM.  The base and perturbed curves nearly overlap within
  each propagation method because the initial perturbation is only
  $5\times10^{-6}$; the following page resolves their separation directly.}
\end{figure}

{\footnotesize\noindent
The uncorrected moment EOM produce bounded but quantitatively different
occupation and energy-partition trajectories from the exact Hamiltonian.
The nearly constant total internal energy does not imply agreement in the
individual observables.\par}

\clearpage

\begin{center}
\small
\textbf{Observable sensitivity to an electronic-imbalance initial perturbation.}
$t_{\rm hop}=1$, $\omega_{\rm ph}=0.5$, $g/t_{\rm hop}=0.6124$
($\lambda=1.5$), $v/t_{\rm hop}=1$; wavefunction perturbation amplitude
$5\times10^{-6}$.
\end{center}

\begin{figure}[htbp]
  \centering
  \includegraphics[width=0.92\textwidth]{output/plots/paper_v_results_progression_20260804/initial_condition_raw_electronic-drive_observable_separation.pdf}
  \caption{Absolute observable separation between the perturbed and base
  initial states.  The black exact-Hamiltonian differences remain near the
  initial $10^{-6}$ scale, while the orange uncorrected-EOM differences grow
  to order $10^{-1}$ in the occupations and energy partitions.  Total
  internal energy remains nearly fixed because both raw trajectories conserve
  their respective initial energies despite separating in the other
  observables.}
\end{figure}

{\footnotesize\noindent
The largest raw observable separations are $0.1094$ in each site occupation,
$0.1748\,t_{\rm hop}$ in electronic energy,
$0.0904\,t_{\rm hop}$ in phonon energy, and
$0.1558\,t_{\rm hop}$ in electron--phonon energy.  The corresponding exact
differences remain below $8.3\times10^{-6}$.\par}

\clearpage

\begin{center}
\small
\textbf{Actual observables for the relative-phonon-position initial perturbation.}
The Hamiltonian, initial ground state, perturbation amplitude, and observable
definitions are identical to the preceding electronic-imbalance test.
\end{center}

\begin{figure}[htbp]
  \centering
  \includegraphics[width=0.92\textwidth]{output/plots/paper_v_results_progression_20260804/initial_condition_raw_relative-phonon-position_observables.pdf}
  \caption{Actual observable trajectories for the base and perturbed initial
  states under exact cutoff-16 Hamiltonian propagation and the uncorrected
  archive moment EOM.  The next page isolates the small within-method
  differences that are difficult to see on the absolute observable scale.}
\end{figure}

\clearpage

\begin{center}
\small
\textbf{Observable sensitivity to a relative-phonon-position initial perturbation.}
The Hamiltonian, initial ground state, perturbation amplitude, and observable
definitions are identical to the preceding page.
\end{center}

\begin{figure}[htbp]
  \centering
  \includegraphics[width=0.92\textwidth]{output/plots/paper_v_results_progression_20260804/initial_condition_raw_relative-phonon-position_observable_separation.pdf}
  \caption{Absolute observable separation for the relative phonon-position
  perturbation.  The exact differences again stay near $10^{-6}$, while the
  uncorrected moment EOM amplify the occupation and energy-partition
  differences over time.}
\end{figure}

{\footnotesize\noindent
The largest raw observable separations are $0.0762$ in each site occupation,
$0.1365\,t_{\rm hop}$ in electronic energy,
$0.0784\,t_{\rm hop}$ in phonon energy, and
$0.1768\,t_{\rm hop}$ in electron--phonon energy.  The corresponding exact
differences remain below $6.7\times10^{-6}$.\par}

\clearpage

\begin{center}
\small
\textbf{Nearby-initial-condition sensitivity of the uncorrected archive EOM.}
$t_{\rm hop}=1$, $\omega_{\rm ph}=0.5$, $g/t_{\rm hop}=0.6124$
($\lambda=1.5$), $v/t_{\rm hop}=1$; cutoff-16 exact-correlated
initialization; matched RK4 step $0.02$ through $t=100$.
\end{center}

{\footnotesize\noindent
From the exact ground state $|\psi_0\rangle$, each nearby physical state is
$|\psi_\epsilon\rangle=(|\psi_0\rangle+\epsilon|\tau_O\rangle)/
\sqrt{1+\epsilon^2}$, where
$|\tau_O\rangle\propto(O-\langle O\rangle)|\psi_0\rangle$ and
$\epsilon=5\times10^{-6}$.  The two generators are the electronic site
imbalance and the relative phonon position.  Each ket is contracted once into
the full 31-coordinate initial state; no individual moment coordinate is
changed by hand.\par}

\begin{figure}[htbp]
  \centering
  \includegraphics[width=0.91\textwidth]{output/plots/paper_v_results_progression_20260804/initial_condition_raw_matched_amplification.pdf}
  \caption{Matched nearby-state test for the uncorrected archive EOM.  Black
  propagates the same physical pairs with the exact cutoff-16 Hamiltonian;
  orange propagates their contracted coordinates with the raw moment EOM.
  Gray denotes the interval after the first sampled retained-constraint
  violation.}
\end{figure}

\begin{figure}[htbp]
  \centering
  \includegraphics[width=0.91\textwidth]{output/plots/paper_v_results_progression_20260804/initial_condition_raw_physicality.pdf}
  \caption{Loss of representability in the raw base trajectory.  The
  correlation-trace equality is violated by the first stored sample at
  $t=0.1\,t_{\rm hop}^{-1}$; the joint Gram is negative at the next sample,
  consistent with the separately refined crossing at
  $t=0.1607\,t_{\rm hop}^{-1}$.}
\end{figure}

\vspace{-3pt}
{\footnotesize\noindent
The raw mathematical ODE amplifies the two initial separations by factors of
$4.39\times10^4$ and $2.55\times10^4$ by $t=100$, whereas the exact contracted
separations remain of order one.  This establishes extreme sensitivity of the
formal uncorrected ODE continuation.  Because the raw trajectory leaves the
representable moment domain before the large amplification develops, it does
not establish physical chaos of the electron--phonon system.\par}

\clearpage

\begin{center}
\small
\textbf{Initial-condition sensitivity of the representability-corrected archive EOM.}
$t_{\rm hop}=1$, $\omega_{\rm ph}=0.5$, $g/t_{\rm hop}=0.6124$
($\lambda=1.5$), $v/t_{\rm hop}=1$; cutoff-16 exact-correlated
initialization and matched physical perturbations.
\end{center}

\begin{figure}[htbp]
  \centering
  \includegraphics[width=0.91\textwidth]{output/plots/paper_v_results_progression_20260804/initial_condition_matched_amplification.pdf}
  \caption{Matched nearby-state test without perturbation rescaling.  The
  same electronic-drive and relative-phonon perturbations are propagated by
  exact cutoff-16 Hamiltonian dynamics and by the corrected 31-coordinate
  moment EOM.  The plotted distance is the combined lifted Frobenius distance
  divided by its initial value.}
\end{figure}

\begin{figure}[htbp]
  \centering
  \includegraphics[width=0.91\textwidth]{output/plots/paper_v_results_progression_20260804/initial_condition_lyapunov_convergence.pdf}
  \caption{Post-pulse Benettin calculation for the corrected moment EOM.
  Repeated tangent-space rescaling through $t=1000\,t_{\rm hop}^{-1}$ gives
  the cumulative exponent (left) and a trailing 250-unit estimate (right).
  The two independent perturbation directions converge to the same positive
  rate.}
\end{figure}

\vspace{-3pt}
{\footnotesize\noindent
Under exact Hamiltonian propagation, the two contracted separations never
exceed factors of $1.385$ and $1.490$ and end below their initial values.  The
corrected moment EOM instead amplify them by factors of $347.2$ and $606.8$ by
$t=100$.  Its cumulative post-pulse exponent is $0.02846\,t_{\rm hop}$, while
the final trailing-window rate remains positive at $0.01891\,t_{\rm hop}$.
The sensitivity therefore belongs to the tested approximate corrected moment
vector field, not to the matched exact dynamics.\par}

\clearpage

\begin{center}
\small
\textbf{Latest result: unconstrained carried higher moments through the complete pulse.}
$t_{\rm hop}=1$, $\omega_{\rm ph}=0.5$, $g/t_{\rm hop}=0.6124$
($\lambda=1.5$), $v/t_{\rm hop}=1$; single Gaussian pulse of width $1$;
exact-correlated cutoff-16 initialization; $h=0.0025$;
$0\le t\,t_{\rm hop}\le4$.
\end{center}

\begin{figure}[htbp]
  \centering
  \includegraphics[width=0.88\textwidth]{output/plots/paper_v_results_progression_20260804/cwrmf_no_guard_t4_combined.pdf}
  \caption{Full-pulse no-guard audit. Panels (a)--(f) compare exact cutoff-16
  Hamiltonian propagation, the raw 31-coordinate archive EOM, and the
  384-coordinate archive-plus-$K/P/D$ higher-moment model with its 62-row
  Gram guard disabled. Panel (g) tracks the extended and retained
  representability margins. Panel (h) shows instantaneous dynamic
  31-coordinate error. Panel (i) estimates how minimum-norm linearized repair
  of each negative extended-Gram mode would change the retained velocity and,
  after a $0.05\,t_{\rm hop}^{-1}$ entrance-channel interval, the $C$ rate.
  This mode-coupling audit is offline and does not alter the trajectory.}
\end{figure}

\begin{center}
  \scriptsize
  \setlength{\tabcolsep}{3.0pt}
  \begin{tabular}{@{}lrrrrrrr@{}}
    \toprule
    Model & 31 RMS & $C$ RMS & $n_0$ RMS & $E_{\rm int}$ RMS
      & $\min\lambda(\mathcal G)$ & $\min\lambda(\mathcal G_{62})$
      & runtime \\
    \midrule
    Exact cutoff-16 & 0 & 0 & 0 & 0 & positive & --- & offline \\
    Raw 31-coordinate archive & $0.11090$ & $0.07660$ & $0.02961$
      & $0.02759$ & $-0.5173$ & --- & not isolated \\
    Carried moments, no Gram guard & $0.48049$ & $0.47234$ & $0.04864$
      & $0.01388$ & $-17.615$ & $-26.445$ & $153.5$ s \\
    \bottomrule
  \end{tabular}
\end{center}
\vspace{-4pt}
{\footnotesize\noindent
The no-guard route remains highly accurate only on the short interval reported
on the next page. Over the complete pulse, the failure sequence is
quantitatively ordered: $\mathcal G_{62}$ crosses the $-10^{-8}$ numerical
floor at the first step; predicted negative-mode feedback into the retained
velocity exceeds $1\%$ at $t=0.10$ and $10\%$ at $t=0.50$; its two-stage
effect on the $C$ rate exceeds $1\%$ at $t=0.55$ and $10\%$ at $t=1.10$;
the instantaneous coordinate RMS exceeds $10^{-3}$ at $t=1.0825$; and the
retained $\mathcal G$ crosses negative at $t=1.355$. Thus extended-cone
violation is initially hidden from the reported observables but later enters
them through the entrance-channel dynamics. Unconstrained higher moments are
an effective fast predictor, but not a complete long-pulse closure.\par}

\clearpage

\begin{center}
\small
\textbf{Latest result: carried higher moments with and without the extended Gram guard.}
$t_{\rm hop}=1$, $\omega_{\rm ph}=0.5$, $g/t_{\rm hop}=0.6124$
($\lambda=1.5$), $v/t_{\rm hop}=1$; single Gaussian pulse of width $1$;
exact-correlated cutoff-16 initialization; matched interval
$0\le t\,t_{\rm hop}\le0.5$.
\end{center}

\begin{figure}[htbp]
  \centering
  \includegraphics[width=0.90\textwidth]{output/plots/paper_v_results_progression_20260804/cwrmf_balanced_t05_observables.pdf}
  \caption{Stored matched trajectories. Black is exact cutoff-16 Hamiltonian
  propagation. Red dashed propagates only the 31 archive coordinates with the
  uncorrected archive EOM. Gray dotted is the strict carried higher-moment
  closure at $h=0.0025$. Blue is the balanced realization of the same closure,
  using $h=0.01$ through $t=0.4$ and $h=0.005$ thereafter. Green dash-dot
  propagates the same retained and higher-moment rates at $h=0.0025$ while
  bypassing the 62-row Gram optimization and endpoint rejection. Exact data
  were used only for offline scoring.}
\end{figure}

\begin{center}
  \tiny
  \setlength{\tabcolsep}{1.8pt}
  \begin{tabular}{@{}lrrrrrrrr@{}}
    \toprule
    Model & step & 31 RMS & $C$ RMS & $n_0$ RMS & $E_{\rm int}$ RMS
      & $\min\lambda(\mathcal G)$ & $\min\lambda(\mathcal G_{62})$
      & solver time \\
    \midrule
    Exact cutoff-16 & adaptive & --- & --- & --- & ---
      & $2.89\times10^{-5}$ & --- & offline \\
    Raw 31-coordinate archive & adaptive & $2.37\times10^{-2}$
      & $3.37\times10^{-2}$ & $4.48\times10^{-4}$
      & $4.89\times10^{-4}$ & $-6.28\times10^{-2}$ & --- & not isolated \\
    Carried moments, no Gram guard & $0.0025$ & $3.19\times10^{-5}$
      & $8.13\times10^{-6}$ & $2.07\times10^{-7}$
      & $1.06\times10^{-7}$ & $-1.03\times10^{-15}$
      & $-1.27\times10^{-2}$ & $18.9$ s \\
    Strict carried moments & $0.0025$ & $3.19\times10^{-5}$
      & $8.14\times10^{-6}$ & $2.06\times10^{-7}$
      & $1.10\times10^{-7}$ & $-1.32\times10^{-15}$
      & $-7.58\times10^{-12}$ & $12.14$ h \\
    Balanced carried moments & $0.01\to0.005$ & $3.56\times10^{-5}$
      & $9.33\times10^{-6}$ & $3.26\times10^{-6}$
      & $2.02\times10^{-6}$ & $-2.20\times10^{-15}$
      & $-9.41\times10^{-9}$ & $4.88$ h \\
    \bottomrule
  \end{tabular}
\end{center}
\vspace{-4pt}
{\footnotesize\noindent
The reported dynamic RMS errors subtract the fixed initial embedding offset.
The retained joint Gram is $\mathcal G$; $\mathcal G_{62}$ is the extended
Gram containing the carried higher moments. The carried routes use the
archive EOM plus the connected electron--two-phonon,
same-spin Pauli-covariance, and opposite-spin covariance contributions
evaluated from the carried higher moments. In the no-guard ablation, every
commutator-readable hidden rate is propagated directly and unresolved rates
retain the minimum-motion value zero. Its retained trajectory differs from
the strict result by only $6.57\times10^{-7}$ coordinate RMS and its half-step
refinement differs by $7.76\times10^{-8}$, while runtime falls by a factor of
$2308$. Thus the cone solve does not create the observed short-time accuracy.
It enforces higher-order representability: without it, $\mathcal G_{62}$ is
already below $-10^{-8}$ at the first step and reaches $-1.27\times10^{-2}$,
even though the retained $\mathcal G$ remains positive to roundoff. The
no-guard trajectory is therefore a fast and accurate short-time diagnostic,
not yet a physical long-horizon closure.\par}

\clearpage

\begin{center}
\small
\textbf{Latest result: common-scale and normalized spectra for four strong-coupling routes.}
$t_{\rm hop}=1$, $\omega_{\rm ph}=0.5$, $g/t_{\rm hop}=0.6124$
($\lambda=1.5$), $v/t_{\rm hop}=1$; single Gaussian pulse of width $1$;
exact-correlated cutoff-16 initialization; common post-pulse window
$4\le t\,t_{\rm hop}\le20$.
\end{center}

\begin{figure}[htbp]
  \centering
  \includegraphics[width=0.91\textwidth]{output/plots/paper_v_results_progression_20260804/archive_m4_four_route_polarization_spectrum.pdf}
  \caption{Electronic-polarization spectra. Panel (a) uses one common
  normalization: the area under the exact Hann-windowed spectrum is one, so
  relative response power is preserved. Panel (b) normalizes every spectrum
  independently to unit area and therefore compares shape only. The four
  lines are
  separate propagated models: exact cutoff-16 Hamiltonian evolution;
  uncorrected 31-coordinate archive EOM; the regular 31-coordinate EOM with
  its minimum-Euclidean-norm joint-Gram velocity correction; and the
  implemented 60-coordinate APCM entrance-layer prototype with positive
  fourth-moment $M_4$ completion, hidden-stage retraction, and the retained
  joint-Gram controller.  Circles mark the strongest peak in
  $1.5\le\omega/t_{\rm hop}\le3.5$.}
\end{figure}

\begin{center}
  \scriptsize
  \setlength{\tabcolsep}{3.5pt}
  \begin{tabular}{@{}lrrrrr@{}}
    \toprule
    Propagated model & $H$ & RMS amplitude/exact
      & $\omega_{\rm pk}/t_{\rm hop}$ & $\Gamma_{\rm obs}/t_{\rm hop}$
      & $W_{\rm band}/W_{\rm band}^{\rm ex}$ \\
    \midrule
    Exact cutoff-16 Hamiltonian & 0 & 1.000 & 2.3411 & 0.4619 & 1.0000 \\
    Archive EOM (uncorrected) & 0.1701 & 1.522 & 2.3549 & 0.5450 & 3.5956 \\
    Regular EOM correction (31D joint-Gram)
      & 0.5868 & 0.873 & 2.6976 & 0.4636 & 0.0269 \\
    McLachlan-type $M_4$ correction (APCM prototype)
      & 0.3708 & 0.868 & 2.4863 & 0.7905 & 0.3600 \\
    \bottomrule
  \end{tabular}
\end{center}
\vspace{-4pt}
{\footnotesize\noindent
$H$ is the Hellinger distance from the complete exact unit-area spectrum;
$W_{\rm band}/W_{\rm band}^{\rm ex}$ is the unnormalized Hann-windowed
power in $1.5\le\omega/t_{\rm hop}\le3.5$ relative to exact.
The common $16\,t_{\rm hop}^{-1}$ window gives
$\Delta\omega=0.3879\,t_{\rm hop}$; consequently
$\Gamma_{\rm obs}$ is a resolution-dominated finite-window FWHM, not a
deconvolved lifetime.  The raw archive EOM give the closest short-window peak
and spectral shape but produce $3.60$ times the exact electronic-band power.
The regular 31D correction's apparently exact width is not an accuracy
success: it leaves only $2.69\%$ of the exact electronic-band power.  The
$M_4$ prototype retains $36.0\%$, but its peak is shifted and broader. This prototype does
not yet implement the full adaptive moment-metric McLachlan projection from
the theoretical APCM design.\par}

\clearpage

\begin{center}
\small
\textbf{Long-window result: interaction-dependent polarization peaks and observed broadening.}
$t_{\rm hop}=1$, $\omega_{\rm ph}=0.5$, $v/t_{\rm hop}=1$;
single Gaussian pulse of width $1$; exact-correlated cutoff-16 ground-state
initialization; common post-pulse window $10\le t\,t_{\rm hop}\le100$.
\end{center}

\begin{figure}[htbp]
  \centering
  \includegraphics[width=0.94\textwidth]{output/plots/paper_v_results_progression_20260804/archive_polarization_peaks_g_comparison.pdf}
  \caption{Power spectrum of the electronic polarization
  $P=(\rho_{11}-\rho_{00})/2$. The left column uses one exact-referenced
  common scale within each coupling regime and preserves relative power; the
  right column normalizes every spectrum independently to unit area and
  compares shape only. All curves use the same
  $0.2\,t_{\rm hop}^{-1}$ sampling, mean subtraction, and Hann window.
  Circles mark the strongest peak in the shaded electronic-frequency band
  $1.5\le\omega/t_{\rm hop}\le3.5$. Exact cutoff-16 data are used only for
  offline scoring.}
\end{figure}

\begin{center}
  \scriptsize
  \setlength{\tabcolsep}{3.6pt}
  \begin{tabular}{@{}llrrrrr@{}}
    \toprule
    $g/t_{\rm hop}$ & Model & $H$ & RMS amplitude/exact
      & $\omega_{\rm pk}/t_{\rm hop}$
      & $\Gamma_{\rm obs}/t_{\rm hop}$
      & $W_{\rm band}/W_{\rm band}^{\rm ex}$ \\
    \midrule
    $0.3536$ & Exact cutoff-16 & 0 & 1.000 & 2.0901 & 0.0940 & 1.0000 \\
             & Uncorrected archive EOM & 0.1464 & 1.086 & 2.0872 & 0.1171 & 1.6884 \\
             & Regular EOM correction (31D joint-Gram) & 0.1958 & 0.644 & 2.0889 & 0.0944 & 0.2760 \\
    \addlinespace
    $0.6124$ & Exact cutoff-16 & 0 & 1.000 & 2.3561 & 0.1145 & 1.0000 \\
             & Uncorrected archive EOM & 0.6381 & 0.750 & 2.9399 & 0.1557 & 3.5834 \\
             & Regular EOM correction (31D joint-Gram) & 0.1782 & 0.746 & 2.4811 & 0.1395 & 0.0223 \\
    \bottomrule
  \end{tabular}
\end{center}
\vspace{-3pt}
{\footnotesize\noindent
$H$ is the Hellinger distance between the complete unit-area spectra;
$W_{\rm band}/W_{\rm band}^{\rm ex}$ is the common-scale electronic-band
power relative to exact.
$\Gamma_{\rm obs}$ is the interpolated full width at half maximum of the
marked peak. It includes the finite-window resolution
$\Delta\omega=0.0697\,t_{\rm hop}$ and Hann-window broadening, so it is an
observed spectral width rather than a deconvolved lifetime. At weak coupling
the archive EOM locate the peak accurately but broaden it by about $25\%$
and produce $1.69$ times the exact band power. The correction retains the
peak location but only $27.6\%$ of the exact band power.
At strong coupling they shift the peak upward by $0.584\,t_{\rm hop}$,
broaden it by about $36\%$, and overpopulate the electronic-frequency band.
The regular 31D EOM correction restores representability and improves the
strong-coupling unit-area spectral shape and peak location, but leaves only
$2.23\%$ of the exact electronic-band power; it therefore does not repair the
missing dynamical closure.\par}

\clearpage

\begin{center}
\small
\textbf{Latest result: valid higher-moment checkpoint at $t=31$.}
$t_{\rm hop}=1$, $\gamma=0.5$, $\lambda=1.5$, $g/t_{\rm hop}=0.6124$;
single Gaussian pulse amplitude $v/t_{\rm hop}=1$, width $1$;
phonon cutoff $16$; 60-variable entrance-layer model with both the retained
joint-Gram and extended $M_4$ cones; SSPRK(3,3) with
$h=0.0025\,t_{\rm hop}^{-1}$.
\end{center}

\begin{figure}[htbp]
  \centering
  \includegraphics[width=0.86\textwidth]{output/plots/paper_v_results_progression_20260804/apcm_t31_checkpoint_observables.pdf}
  \caption{Interrupted long-horizon calculation. Black is the stored exact
  cutoff-16 Hamiltonian reference through $t=31$; red dash-dot is the stored
  uncorrected archive EOM; blue dashed is the autonomous both-cones trajectory
  through $t=20$; and the blue diamond is its finite, resumable
  checkpoint at $t=31$. No line is drawn over
  $20<t\,t_{\rm hop}<31$ because the interrupted process checkpointed its
  state and fourth-moment frontier but did not serialize intermediate plotting
  samples. At $t=31$, the retained joint-Gram and $M_4$ margins were
  $1.224\times10^{-4}$ and $3.032\times10^{-5}$, respectively.}
\end{figure}

\begin{center}
  \small
  \setlength{\tabcolsep}{5pt}
  \begin{tabular}{@{}lrrrr@{}}
    \toprule
    Model & 31-coordinate RMS & $C$ Frobenius error & $|\Delta n_0|$
      & $|\Delta E_{\rm int}|/t_{\rm hop}$ \\
    \midrule
    Uncorrected archive
      & $0.43993$ & $0.45749$ & $1.05629$ & $0.015011$ \\
    Both-cones checkpoint
      & $0.29204$ & $0.27406$ & $0.66563$ & $9.52\times10^{-4}$ \\
    \bottomrule
  \end{tabular}
\end{center}
\vspace{-3pt}
{\footnotesize\noindent
Endpoint errors are measured against the exact cutoff-16 state at $t=31$.
The both-cones checkpoint remains representable, whereas the uncorrected
archive joint-Gram eigenvalue is $-0.48988$. The correction improves the
aggregate coordinates, $C$, occupation, and total energy, but substantial
dynamical error remains and some individual energy-component errors increase.
The next run resumes from this exact stored state rather than repeating
$0\le t\,t_{\rm hop}\le31$.\par}

\clearpage

\begin{center}
\small
\textbf{Latest result: accelerated higher-moment cone routes through $t=20$.}
$t_{\rm hop}=1$, $\gamma=0.5$, $\lambda=1.5$, $g/t_{\rm hop}=0.6124$;
single Gaussian pulse amplitude $v/t_{\rm hop}=1$, width $1$;
phonon cutoff $16$; 60-variable entrance-layer model;
SSPRK(3,3) with $h=0.0025\,t_{\rm hop}^{-1}$ and exact spin-exchange
block decomposition of the $M_4$ cone.
\end{center}

\begin{figure}[htbp]
  \centering
  \includegraphics[width=0.82\textwidth]{output/plots/paper_v_results_progression_20260804/apcm_t20_observables.pdf}
  \par\vspace{-3pt}
  \includegraphics[width=0.74\textwidth]{output/plots/paper_v_results_progression_20260804/apcm_t20_diagnostics.pdf}
  \caption{Matched autonomous trajectories and diagnostics from the same
  exact-correlated initial contraction. Solid black is exact cutoff-16
  Hamiltonian propagation, used only after the two rollouts for scoring. Red
  dash-dot is the stored uncorrected archive EOM. Blue dashed is the
  higher-moment model with both the retained joint-Gram and extended $M_4$
  cones; green dotted uses only the $M_4$ cone. Panels (g)--(l)
  show instantaneous errors, both cone margins, retained-cone velocity
  correction, and hidden-stage retraction. Small negative $M_4$ values remain
  within the declared $10^{-6}$ conic backward-error tolerance.}
\end{figure}

\begin{center}
  \tiny
  \setlength{\tabcolsep}{2.1pt}
  \begin{tabular}{@{}lrrrrrrr@{}}
    \toprule
    Model & 31 scalar RMS & $C$ RMS & $n_0$ RMS & $E_{\rm int}$ RMS
      & $\min\lambda(\mathcal G)$ & $\min\lambda(M_4)$ & runtime \\
    \midrule
    Exact cutoff-16 Hamiltonian
      & 0 & 0 & 0 & 0 & $2.85\times10^{-5}$ & -- & 5.2 s \\
    Uncorrected archive EOM
      & 0.19320 & 0.39095 & 0.28617 & 0.018236
      & $-1.0281$ & -- & not isolated \\
    Higher moment + both cones
      & 0.12861 & 0.14774 & 0.20469 & 0.000940
      & $1.37\times10^{-5}$ & $-9.89\times10^{-7}$ & 5832 s \\
    Higher moment + $M_4$ cone only
      & 0.12884 & 0.14602 & 0.19201 & 0.000940
      & $2.67\times10^{-5}$ & $-9.83\times10^{-7}$ & 5955 s \\
    \bottomrule
  \end{tabular}
\end{center}
\vspace{-4pt}
{\footnotesize\noindent
Time-RMS errors over $0\le t\,t_{\rm hop}\le20$. The both-cones route is
slightly better in aggregate coordinate error; the $M_4$-only route is
slightly better in occupation and $C$. Neither trajectory crosses the
amplitude threshold or loses retained joint-Gram positivity. Maximum
retained-cone correction and hidden-stage retraction norms were $0.0702$ and
$0.00147$ with both cones, versus $0$ and $0.00123$ with $M_4$ only.\par}

\clearpage

\begin{center}
\small
\textbf{Latest result: matched retained- and extended-cone ablation through $t=4$.}
$t_{\rm hop}=1$, $\gamma=0.5$, $\lambda=1.5$, $g/t_{\rm hop}=0.6124$;
single Gaussian pulse amplitude $v/t_{\rm hop}=1$, width $1$;
phonon cutoff $16$; 60-variable entrance-layer model;
SSPRK(3,3) with $h=0.0025\,t_{\rm hop}^{-1}$.
\end{center}

\begin{figure}[htbp]
  \centering
  \includegraphics[width=0.93\textwidth]{output/plots/paper_v_results_progression_20260804/apcm_t4_matched_observables.pdf}
  \caption{Matched trajectories from the same exact-correlated initial
  contraction and single-pulse protocol.  The comparison contains exact
  cutoff-16 Hamiltonian propagation, the raw 31-coordinate archive EOM, that
  EOM with its retained joint-Gram cone controller, and the higher-moment
  model under all four combinations of the retained joint-Gram cone and the
  extended $M_4$ cone.  Exact data enter only after each autonomous rollout.
  The stored 31-coordinate curves use $h=0.01$ and are interpolated onto the
  $h=0.0025$ display grid; their errors are evaluated on their native grid.}
\end{figure}

\begin{table}[htbp]
  \centering
  \tiny
  \setlength{\tabcolsep}{2.2pt}
  \begin{tabular}{@{}lrrrrrrr@{}}
    \toprule
    Model & 31 scalar RMS & $C$ RMS & $n_0$ RMS & $E_{\rm int}$ RMS
        & $\min\lambda(\mathcal G)$ & $\min\lambda(M_4)$ & runtime \\
    \midrule
    Exact cutoff-16 Hamiltonian
      & 0 & 0 & 0 & 0
      & $2.85\times10^{-5}$ & -- & 1.0--1.5 s \\
    31-coordinate raw EOM
      & 0.11085 & 0.23483 & 0.05924 & 0.02759
      & $-0.5173$ & -- & not isolated$^{*}$ \\
    31-coordinate + retained cone
      & 0.11172 & 0.19380 & 0.11494 & 0.05123
      & $2.06\times10^{-5}$ & -- & not isolated$^{*}$ \\
    Higher moment: no cones
      & 0.04995 & 0.10663 & 0.04806 & 0.02509
      & $-8.23\times10^{-3}$ & $-1.487$ & 27.3 s \\
    Higher moment + retained cone
      & 0.04647 & 0.10323 & 0.05179 & 0.02697
      & $1.03\times10^{-5}$ & $-1.473$ & 52.1 s \\
    Higher moment + $M_4$ cone
      & 0.01258 & 0.04743 & 0.01512 & 0.000821
      & $2.89\times10^{-5}$ & $-9.68\times10^{-7}$ & 3268.6 s \\
    Higher moment + both cones
      & 0.01258 & 0.04725 & 0.01385 & 0.000796
      & $2.89\times10^{-5}$ & $-8.76\times10^{-7}$ & 3399.6 s \\
    \bottomrule
  \end{tabular}
  \caption{Errors are time-RMS values over $0\le t\,t_{\rm hop}\le4$;
  $C$ uses the lifted Frobenius norm.  The small negative $M_4$ minima of the
  two positive-completion runs remain within the declared $10^{-6}$ conic
  backward-error tolerance.  A dash means that $M_4$ is not a propagated
  object in that model. $^{*}$The stored 31-coordinate artifact records
  40.7 s for the exact scorer and four EOM lanes together, not separate
  per-lane runtimes.}
\end{table}

\clearpage

\begin{center}
\small
\textbf{Latest result: adaptive packet continuation through $t=98$.}
$t_{\rm hop}=1$, $\gamma=0.5$, $\lambda=1.5$, $g/t_{\rm hop}=0.6124$;
double-pulse amplitude $v/t_{\rm hop}=1$, delays $0$ and $8$, width $1$;
phonon cutoff $16$.
\end{center}

\begin{figure}[htbp]
  \centering
  \includegraphics[width=0.82\textwidth]{output/plots/paper_v_results_progression_20260804/adaptive_t98_observables.pdf}
  \caption{Contracted observables through the last complete checkpoint at
  $t=98$.  The trajectory resumed from the stored $t=40$, $K=11$ state and
  reached $K=27$ coherent packets per electronic branch under the unchanged,
  uncapped admission rule. The matched raw archive EOM reaches the declared
  amplitude threshold at $t=90.07$. Exact cutoff-16 data enter only after
  autonomous packet-ket propagation for scoring.}
\end{figure}

\begin{table}[htbp]
  \centering
  \scriptsize
  \begin{tabular}{lrrrrrrrr}
    \toprule
    Model & final $K$ & $\bar K$ & 31 RMS & $C$ RMS & $n_0$ RMS & $E_{\rm ph}$ RMS & $E_{\rm e-ph}$ RMS & $E_{\rm int}$ RMS \\
    \midrule
    Ordinary packet & 6 & 5.93 & 0.14079 & 0.14744 & 0.04522 & 0.06101 & 0.08840 & 0.001188 \\
    Adaptive packet & 27 & 12.57 & 0.06365 & 0.05653 & 0.01360 & 0.02835 & 0.03379 & 0.000579 \\
    \bottomrule
  \end{tabular}
  \caption{Time-RMS errors over $0\le t\,t_{\rm hop}\le98$ against exact
  cutoff-16 Hamiltonian propagation.  The adaptive ket remained representable
  without the archive-EOM minimum-norm physicality correction; its minimum
  exact-state fidelity was $0.92368$.  The rapid late growth to $K=27$ also
  indicates that the frozen residual trigger is aggressive at this horizon.}
\end{table}

\clearpage

\begin{center}
\small
\textbf{Latest result: uncapped adaptive packet growth through $t=40$.}
$t_{\rm hop}=1$, $\gamma=0.5$, $\lambda=1.5$, $g/t_{\rm hop}=0.6124$;
double-pulse amplitude $v/t_{\rm hop}=1$, delays $0$ and $8$, width $1$;
phonon cutoff $16$.
\end{center}

\begin{figure}[htbp]
  \centering
  \includegraphics[width=0.82\textwidth]{output/plots/paper_v_results_progression_20260804/adaptive_t40_observables.pdf}
  \caption{Observable trajectories through $t=40$.  The frozen gate began at
  $K=6$, admitted state-continuously at $t=8,19,22,26.5,38$, and ended at
  $K=11$ without a packet ceiling.  The fixed $K_{\max}=10$ trajectory is the
  strongest aggregate fixed-capacity comparator; the dashed brown curve is
  the matched raw archive EOM. Exact cutoff-16 data enter only after
  autonomous propagation for scoring.}
\end{figure}

\begin{table}[htbp]
  \centering
  \scriptsize
  \begin{tabular}{lrrrrrrrr}
    \toprule
    Model & final $K$ & $\bar K$ & 31 RMS & $C$ RMS & $n_0$ RMS & $E_{\rm ph}$ RMS & $E_{\rm e-ph}$ RMS & $E_{\rm int}$ RMS \\
    \midrule
    Ordinary $K_{\max}=6$ & 6 & 5.83 & 0.05338 & 0.06388 & 0.01201 & 0.01522 & 0.03296 & 0.000742 \\
    Fixed $K_{\max}=8$ & 8 & 7.34 & 0.05069 & 0.04556 & 0.01030 & 0.02934 & 0.03489 & 0.000757 \\
    Fixed $K_{\max}=10$ & 10 & 8.49 & 0.03907 & 0.02769 & 0.00606 & 0.03099 & 0.03189 & 0.000760 \\
    Fixed $K_{\max}=12$ & 12 & 9.31 & 0.04859 & 0.03859 & 0.00745 & 0.03079 & 0.03400 & 0.000799 \\
    Adaptive gate & 11 & 8.17 & 0.03763 & 0.02885 & 0.00501 & 0.02447 & 0.02672 & 0.000654 \\
    \bottomrule
  \end{tabular}
  \caption{Time-RMS errors over $0\le t\,t_{\rm hop}\le40$ and time-average
  packets per branch $\bar K$.  The adaptive path gives the lowest aggregate,
  occupation, electron--phonon-energy, and total-energy errors at lower
  average capacity than fixed $K_{\max}=10$; fixed $K_{\max}=10$ retains
  slightly lower $C$-block and electronic-energy error.}
\end{table}

\clearpage

\begin{center}
\small
\textbf{Transfer result: frozen adaptive gate across preparation and drive.}
$t_{\rm hop}=1$, $\gamma=0.5$, $\lambda=1.5$, $g/t_{\rm hop}=0.6124$;
pulse amplitude $v/t_{\rm hop}=1$, width $1$, phonon cutoff $16$, and
$0\le t\,t_{\rm hop}\le20$.
\end{center}

\begin{figure}[htbp]
  \centering
  \includegraphics[width=0.82\textwidth]{output/plots/paper_v_results_progression_20260804/adaptive_transfer_observables.pdf}
  \caption{The unchanged current-state gate on three transfer cells.  Plus
  and minus are nearby physical preparations under the double pulse at
  $t=0$ and $8$; the final row uses the central preparation under only the
  pulse at $t=0$. The gate chose different admission times in each cell, and
  each row includes the raw archive EOM from that row's initial contraction
  and pulse. Exact cutoff-16 contractions enter only after each autonomous
  rollout for scoring.}
\end{figure}

\begin{table}[htbp]
  \centering
  \scriptsize
  \begin{tabular}{llrrrrrr}
    \toprule
    Cell & Model & final $K$ & 31 RMS & $C$ RMS & $n_0$ RMS & $E_{\rm ph}$ RMS & $E_{\rm e-ph}$ RMS \\
    \midrule
    Plus, double & parent & 6 & 0.02662 & 0.02777 & 0.00472 & 0.01723 & 0.02173 \\
                 & adaptive & 8 & 0.02034 & 0.01308 & 0.00143 & 0.01143 & 0.01356 \\
    Minus, double & parent & 6 & 0.03073 & 0.02391 & 0.00390 & 0.01533 & 0.01809 \\
                  & adaptive & 9 & 0.01790 & 0.01466 & 0.00185 & 0.00790 & 0.01117 \\
    Central, single & parent & 6 & 0.01747 & 0.01672 & 0.00253 & 0.01100 & 0.01332 \\
                    & adaptive & 9 & 0.01073 & 0.00976 & 0.00161 & 0.00874 & 0.00984 \\
    \bottomrule
  \end{tabular}
  \caption{Time-RMS errors against exact cutoff-16 Hamiltonian propagation.
  The frozen gate improves aggregate moments, the correlation block, and all
  displayed observables in all three transfer cells.  Its post-initial
  admissions were $(11,12)$ for plus, $(13.5,16.5,18.5)$ for minus, and
  $(8,11,18)$ for the single-pulse cell.}
\end{table}

\clearpage

\begin{center}
\small
\textbf{Development cell: state-adapted mixed-guided packet admission.}
$t_{\rm hop}=1$, $\omega_{\rm ph}/t_{\rm hop}=\gamma=0.5$,
$\lambda=1.5$, $g/t_{\rm hop}=0.6124$;
double-pulse amplitude $v/t_{\rm hop}=1$, delays $0$ and $8$, pulse width $1$;
phonon cutoff $16$ and $0\le t\,t_{\rm hop}\le20$.
\end{center}

\begin{figure}[htbp]
  \centering
  \includegraphics[width=0.98\textwidth]{output/plots/paper_v_results_progression_20260804/mixed_guided_readmission_observables.pdf}
  \caption{Observable trajectories for autonomous packet propagation.  The
  ordinary parent reaches $K=6$ by native-residual admission.  The guided
  models instead add two zero-weight packets at $t=0$ using the current mixed
  electron--phonon residual; the $K=7$ and $K=8$ models add one or two more at
  $t=8$.  The autonomous gate checks the current residual every $0.5$ units
  and admits at $t=8$ and $19$. Every admission preserves the represented
  state exactly; the dashed brown curve is the matched raw archive EOM. Exact
  cutoff-16 data enter only after propagation for scoring.}
\end{figure}

\begin{table}[htbp]
  \centering
  \small
  \begin{tabular}{lrrrrr}
    \toprule
    Model & 31-coordinate RMS & $C$-block RMS & $n_0$ RMS & $E_{\rm ph}$ RMS & $E_{\rm e-ph}$ RMS \\
    \midrule
    Ordinary $K=6$ & 0.02784 & 0.02163 & 0.00271 & 0.01541 & 0.01865 \\
    Guided $K=6$ at $t=0$ & 0.02241 & 0.02442 & 0.00495 & 0.01221 & 0.01706 \\
    Readmit $K=7$ at $t=8$ & 0.01909 & 0.01330 & 0.00135 & 0.00946 & 0.01196 \\
    Readmit $K=8$ at $t=8$ & 0.01697 & 0.01102 & 0.00172 & 0.00952 & 0.01178 \\
    Adaptive $K=6\to7\to8$ & 0.01809 & 0.01291 & 0.00148 & 0.00967 & 0.01206 \\
    \bottomrule
  \end{tabular}
  \caption{Time-RMS errors over $0\le t\,t_{\rm hop}\le20$.  The
  state-adapted readmission restores the post-second-pulse advantage lost by
  the $t=0$-only basis.  The scheduled and autonomous models form a mixed accuracy--cost tradeoff:
  $K=8$ has lower aggregate moment error, while $K=7$ is better for several
  displayed observables and final fidelity.}
\end{table}

\clearpage

\begin{center}
\small
\textbf{Prior result: finite-horizon reciprocal-memory order curve.}
$t_{\rm hop}=1$, $\omega_{\rm ph}/t_{\rm hop}=\gamma=0.5$,
$\lambda=1.5$, $g/t_{\rm hop}=0.6124$;
single-pulse amplitude $v/t_{\rm hop}=1$, delay $0$, pulse width $1$;
phonon cutoff $16$, $\Delta t=0.01\,t_{\rm hop}^{-1}$, and
$0\le t\,t_{\rm hop}\le4$.
\end{center}

\begin{figure}[htbp]
  \centering
  \includegraphics[width=0.98\textwidth]{output/plots/paper_v_results_progression_20260804/finite_horizon_auxiliary_observables.pdf}
  \caption{Observable trajectories for the fixed reciprocal-memory route.
  The $r=78$ model has the lowest coordinate and energy RMS errors among all
  32 distinct finite-horizon union orders $r=20,22,\ldots,80,81$; $r=81$ is
  the complete depth-2 envelope.  The exact cutoff-16 trajectory entered only
  after each autonomous rollout for scoring.}
\end{figure}

\begin{table}[htbp]
  \centering
  \small
  \begin{tabular}{lrrrr}
    \toprule
    Model & Hidden order & Coordinate RMS & Energy RMS & $\min_t\lambda_{\min}(\mathcal G)$ \\
    \midrule
    Raw archive EOM & 0 & 0.110660 & 0.027449 & $-0.51719$ \\
    Pauli-repaired archive EOM & 0 & 0.057742 & 0.096327 & $-0.39979$ \\
    Best finite-horizon union & 78 & 0.347288 & 0.010080 & $-4.50875$ \\
    Complete reciprocal envelope & 81 & 0.362366 & 0.011027 & $-5.06978$ \\
    \bottomrule
  \end{tabular}
  \caption{Sampled RMS errors over $0\le t\,t_{\rm hop}\le4$ and the
  minimum joint-Gram eigenvalue.  The auxiliary route improves total-energy
  agreement but worsens the retained-moment trajectory and representability,
  so rank selection inside this fixed envelope does not repair the closure.}
\end{table}

\clearpage

\begin{center}
\small
\textbf{Hamiltonian and protocol:}
$t_{\rm hop}=1$, $\omega_{\rm ph}/t_{\rm hop}=\gamma=0.5$,
$\lambda=1.5$, $g/t_{\rm hop}=0.6124$
\(\lambda=2g^2/(t_{\rm hop}\omega_{\rm ph})\);
double-pulse amplitude $v/t_{\rm hop}=1$, delays $0$ and $8$, pulse width $1$;
phonon cutoff $16$ and $0\le t\,t_{\rm hop}\le40$.
\end{center}

\begin{figure}[htbp]
  \centering
  \includegraphics[width=0.98\textwidth]{output/plots/paper_v_results_progression_20260804/packet_capacity_observables.pdf}
  \caption{Latest packet-capacity comparison in the strong regime.  Every
  packet trajectory starts from the same fitted four-packet state and uses the
  same double pulse; only the maximum number of coherent packets per electronic
  branch changes.  The black curve is cutoff-16 DOP853 wavefunction
  propagation, and the dashed brown curve is the raw archive EOM from the same
  initial contraction. Shading marks the displayed four-unit pulse windows.}
\end{figure}

\begin{table}[htbp]
  \centering
  \small
  \begin{tabular}{c@{\qquad}rrrrrr}
    \toprule
    $K_{\max}$ & $n_0$ & $E_{\rm e}$ & $E_{\rm ph}$ & $E_{\rm e-ph}$ & $E_{\rm int}$ & source NRMS \\
    \midrule
     6 & 0.01201 & 0.02797 & 0.01522 & 0.03296 & 0.000742 & 0.20427 \\
     8 & 0.01030 & 0.01871 & 0.02934 & 0.03489 & 0.000757 & 0.16053 \\
    10 & 0.00606 & 0.00979 & 0.03099 & 0.03189 & 0.000760 & 0.12695 \\
    12 & 0.00745 & 0.01354 & 0.03079 & 0.03400 & 0.000799 & 0.15172 \\
    \bottomrule
  \end{tabular}
  \caption{Time-RMS observable errors over $0\le t\le40$ relative to the
  cutoff-16 DOP853 trajectory, followed by the normalized RMS error of the
  packet-derived missing correlation source.  Energies are in hopping units.}
\end{table}

\end{document}
